\section{正则化从 Ridge 到 Lasso}
\label{sec:13-Machine-Learning/02-Linear-Models/02}

\subsection{两个几乎相同的温度计，系数却在抵消}

一个化工反应釜上装了两个温度传感器，安装在反应釜外壁相距不到 30 厘米的位置。操作员记录了 200 个时刻的两个传感器读数以及对应的反应速率。用最小二乘回归反应速率对两个温度读数：

\[
\widehat y=8.21x_1-7.34x_2+0.15.
\]

两个系数几乎互为相反数。从物理上看，两个传感器的读数应该高度一致——它们测量的是同一个反应釜、同一时刻的温度。回归系数的含义似乎是：``传感器 1 每升高 1 度，反应速率提高 8.21；但只要传感器 2 同时升高 1 度，这个效果就被抵消掉 7.34。''

这没有物理意义。系数拆分是由数据的偶然波动和传感器的微小差异共同制造的，不是由真实的物理机制决定的。更麻烦的在后头：第二天用同一组传感器采集新一批数据，同样的回归给出

\[
\widehat y=-3.89x_1+4.76x_2+0.22.
\]

系数的符号翻了过来。从 \(8.21\) 到 \(-3.89\)，从 \(-7.34\) 到 \(4.76\)——变动幅度远超任何合理的标准误。

这不是最小二乘出了 bug。两个传感器读数的相关系数高达 0.997，矩阵 \(X^{\trans} X\) 有一个极小的特征值，对应方向是 \(x_1-x_2\approx 0\)。在这个方向上，数据几乎没有提供任何信息去区分两个传感器各自的贡献。最小二乘忠实地求解了正规方程，但解出来的系数对数据中的微小扰动极其敏感——训练数据中的一点点采样波动，就会让系数在两个传感器之间大幅转移。

问题不是算法选错了，是数据没有能力回答``每个传感器各自贡献了多少''这个问题。如果坚持要一个答案，就必须自己补充一个原则：在数据给不出区分的地方，我偏好什么样的解？

\subsection{Ridge：不让弱信息方向的噪声被无限放大}

Ridge 回归在最小二乘的目标上加一项对系数平方和的惩罚：
\[
\widehat\beta_\lambda
=\argmin_\beta
\left\{
\frac12\|y-X\beta\|_2^2+
\frac{\lambda}{2}\|\beta\|_2^2
\right\},
\qquad \lambda\ge 0.
\]

第一项要求拟合数据，第二项禁止系数过大。当 \(\lambda=0\) 时退化为普通最小二乘。一阶条件：
\[
(X^{\trans} X+\lambda I)\widehat\beta_\lambda=X^{\trans} y.
\]

只要 \(\lambda>0\)，矩阵 \(X^{\trans} X+\lambda I\) 就一定可逆——即使 \(X\) 不满列秩，目标在每一个坐标方向上都变成了强凸的。这意味着解是唯一的：不再有多组不同系数叠出相同预测面的问题。

SVD 最清楚地揭示 Ridge 到底做了什么。令 \(X=U\Sigma V^{\trans}\)，其中 \(\Sigma\) 的对角元素是奇异值 \(\sigma_1\ge\sigma_2\ge\cdots\ge\sigma_p\ge 0\)。普通最小二乘在 \(V\) 的旋转坐标系中，沿第 \(j\) 个方向的分量是
\[
\frac{u_j^{\trans} y}{\sigma_j},
\]
其中 \(u_j\) 是 \(U\) 的第 \(j\) 列。当 \(\sigma_j\) 很小时，除以 \(\sigma_j\) 等于用一个大乘数放大 \(u_j^{\trans} y\) 中本来就很小的投影——那就是噪声被放大的数学机制。

Ridge 做了什么？沿着完全相同的方向，Ridge 的解在 \(V\) 的旋转坐标系中写为
\[
\widehat\beta_\lambda
=V\diag
\left(
\frac{\sigma_j}{\sigma_j^2+\lambda}
\right)U^{\trans} y.
\]

比较两个缩放因子。最小二乘的因子是 \(1/\sigma_j\)——对弱方向（\(\sigma_j\) 小）发散。Ridge 的因子是 \(\sigma_j/(\sigma_j^2+\lambda)\)——当 \(\sigma_j\) 很大时约等于 \(1/\sigma_j\)（强方向几乎不受影响）；当 \(\sigma_j\to 0\) 时约等于 \(\sigma_j/\lambda\to 0\)（弱方向被彻底压制）。

回到两个温度传感器。若两个传感器的公共方向（平均值）有大的奇异值，Ridge 几乎完整保留它；两个传感器的差异方向（差值）有极小的奇异值，Ridge 会把它缩到接近零。结果是 Ridge 倾向于把权重平滑分配给两个传感器，而不是让它们互相抵消产生一个巨大的差项。单个系数不再是无偏的（它们被缩向零），但整体预测的方差大幅下降。在未见数据上，更小的方差通常足以弥补引入的偏差——这就是偏差-方差权衡的具体运算。

如果把系数配上一个零均值各向同性 Gaussian 先验，Ridge 的解恰好是该先验下的后验众数（MAP）。这个解释优雅，但依赖于尺度的选择：把温度单位从摄氏度换成开尔文，系数的数值会变化，同一个 \(\lambda\) 就不再代表同一个先验。因此除非原始单位本身承载有意的先验信息，通常应在每个训练折内部将特征标准化到单位方差，并且不惩罚截距——截距只是响应均值在不同特征原点的平移，没有``过大''需要惩罚的问题。

\subsection{Lasso：允许某些坐标恰好等于零}

如果不仅想收缩系数，还希望一部分系数精确为零——直接剔除那些变量，得到一个更紧凑的模型——可以把二次惩罚换成 \(\ell_1\) 惩罚：
\[
\widehat\beta_\lambda
=\argmin_\beta
\left\{
\frac12\|y-X\beta\|_2^2+\lambda\|\beta\|_1
\right\}.
\]

\(\abs{\beta_j}\) 在 \(\beta_j=0\) 处不可导。这恰恰是 Lasso 能产生精确零系数的原因：最优解恰好卡在那个不可导的尖角上。

一维情形最能说明机制。考虑
\[
\min_\beta\ \frac12(\beta-z)^2+\lambda\abs{\beta}.
\]
\(z\) 可以理解为普通最小二乘在那个坐标上给出的单变量估计。

若 \(\beta>0\)，目标可导，一阶条件给出 \(\beta=z-\lambda\)。但这个解必须在 \(\beta>0\) 的前提下自洽，即要求 \(z>\lambda\)。若 \(\beta<0\)，对称地得到 \(\beta=z+\lambda\)，要求 \(z<-\lambda\)。若 \(\abs{z}\le\lambda\)，\(\beta=0\) 落入目标在该点的次微分：
\[
[-z-\lambda,\,-z+\lambda]\ni 0,
\]
因此 \(\beta=0\) 是最优的。合并起来就是软阈值算子：
\[
S_\lambda(z)=\sign(z)(\abs{z}-\lambda)_+.
\]

在正交设计（\(X^{\trans} X=I\)）下，每个坐标的 Lasso 系数独立地等于对其单变量最小二乘估计做软阈值。\(z\) 的绝对值小于 \(\lambda\) 的坐标被精确压到零；超过阈值的坐标被向零方向缩 \(\lambda\)。

几何上，\(\ell_1\) 约束区域 \(\{\beta:\norm{\beta}_1\le t\}\) 是一个在各个坐标轴上有尖角的多面体。损失函数的椭圆等高线更容易在尖角处而不是在平坦面上碰到约束边界——尖角对应某些坐标恰好为零的解。\(\ell_2\) 约束区域是一个光滑的球，等高线在球面上的切点几乎不可能恰好落在坐标轴上——这解释了为什么 Ridge 收缩但不产生精确零。

但``系数为零''有一个容易被过度解读的局限。若两个特征高度相关，Lasso 可能只随机选中其中一个，把另一个压到零。换一批样本，入选者和落选者可能互换。Lasso 选出的是一个对当前数据有效的稀疏预测模型，而不是在宣称``只有这些变量在物理上起作用''。稀疏预测和因果发现之间的鸿沟，不是一个惩罚函数能跨越的。

\subsection{怎样在 Ridge 和 Lasso 之间选择}

Ridge 适合的情形：许多变量都贡献少量信号，相关变量应当共同保留（如相邻的传感器、同一概念的不同测量、图像的相邻像素）。Lasso 适合的情形：相信真实表示在当前坐标系中近似稀疏，且希望得到一个紧凑的、容易解释的模型（如基因表达数据中只有少数基因与表型相关）。

如果既希望成组保留相关变量，又希望稀疏，可以用 Elastic Net：
\[
\Omega(\theta)=\alpha\norm{\theta}_1+(1-\alpha)\norm{\theta}_2^2.
\]
它兼具 L1 的稀疏性和 L2 的群组平滑性，在特征数量远超样本数的情形下通常比单纯的 Lasso 更稳定。

选择 \(\lambda\) 时，必须把数据驱动的所有预处理放在每个训练折内部。常见的错误流程是：先在全部数据上标准化、筛掉低方差特征，再做交叉验证选 \(\lambda\)。此时标准化参数和特征集合已经见到了验证折的标签——\(\lambda\) 的选择被乐观化了。正确流程：在每个训练折内部独立做标准化、特征筛选和 \(\lambda\) 搜索，只在最终的验证折（或外层测试折）上评估。

除了看验证误差最低点的 \(\lambda\)，还应该看整条正则化路径上的行为：

\begin{itemize}
\tightlist
\item
  训练误差和验证误差随 \(\lambda\) 增大如何变化——验证误差通常先降后升，下降段表示正则化在压制过拟合，上升段表示偏差开始超过方差收益。
\item
  Ridge 的系数是否平滑收缩——如果某个系数不随 \(\lambda\) 单调变化，可能特征之间有高度共线性。
\item
  Lasso 的非零系数集合在相邻折和 bootstrap 样本中是否稳定——入选变量反复变化，说明数据不足以稳定地选出特定变量，此时 Lasso 的稀疏性不应被解释为``这些就是重要变量''。
\item
  选择的模型在真正部署分布上是否仍然合理——交叉验证优化的只是训练分布上的期望表现，不保证迁移。
\end{itemize}

数值求解方面：Ridge 在 SVD 分解后可用 \(O(np^2)\) 代价对所有 \(\lambda\) 同时求解。Lasso 常用坐标下降（每次只优化一个坐标，保持其他坐标固定）或近端梯度法。对于 Lasso，目标值变化很小不等于已经收敛到足够的精度——非光滑面上次梯度可以很小，但解仍然离最优点有一段距离。KKT 条件或对偶间隙是更可靠的收敛判据。

\subsection{正则化不是免费的午餐}

正则化用一个偏好换取了稳定性，但偏好本身可能不合适。L2 假设所有方向应该被均匀收缩——如果有些方向确实应该大，这种均匀收缩会引入不必要的偏差。L1 假设表示是稀疏的——如果真实机制中所有变量都有微小贡献，L1 强制稀疏会丢掉累积起来可能很重要的弱信号。

正则化强度 \(\lambda\) 的选择还引入了另一层方差：在有限数据上交叉验证选出的 \(\lambda\)，本身也是一个随机量。样本量较小时，选出的 \(\lambda\) 可能离总体最优值相当远。这是为什么在部署前应该用最终选定的 \(\lambda\) 在全部训练数据上重新拟合，并用从未参与任何选择（包括 \(\lambda\) 的选择）的测试集做最终评估。

最终，正则化不能替代对数据生成机制的理解。两个温度传感器的正确处理可能不是选 Ridge 或 Lasso 去拆开它们——而是取平均值作为单一特征，因为你已经知道它们测量的是同一个物理量。数学工具给出可计算的偏好，但偏好应当表达对问题的认识，而不只是对噪声的防御。
